% =============================================================================
% Part II -- The Boundary  (the climax)
% =============================================================================
\noindent Part~II answers Clause~2 in the only form the present theory can
honestly support: \textbf{where does the construction halt, and what conditional
non-forcing statement can be proved there?} It is the monograph's climax, and
the section where the program's honesty is most directly on the line, because
the quantity at stake is the one whose numerical match ($x_+ =
137.036\ldots$, deferred to Part~III) motivates the whole enterprise. The
temptation to call that match a derivation is exactly what this Part refuses.
What it delivers instead is a \emph{mapped wall}: a route-invariant boundary
around the electromagnetic coupling $\alpha$, built from theorem-grade masonry
around one precisely-located, honestly-marked \etag{opn} gap.

\medskip
\noindent State the shape of the result before its content. The headline\,---\,$\Ftheory$
(the five postulates plus the spine) does not force $\alpha$\,---\,is not a single
theorem. It is a route-invariant boundary assembled from three theorem-grade
components plus one irreducible \etag{opn} kernel; the overall no-go is a
\etag{smc}, and the obstruction it sharpens (the framework's central
foundational obstruction) remains a \etag{fo}.\margetag{fo} The honest
one-line summary is the project's own: \textbf{$\alpha$ is dynamical, not
structural.} It is never ``$\alpha$ is derived,'' and it is never ``$\alpha$ is
proven underivable'' in any absolute sense; \secref{sec:dynamical} states exactly
what it is.

\section{The readout problem}\label{sec:readout}
Part~I built the master quadratic $\masterquad$ as a pure object of the algebraic
spine (\secref{sec:masterquad}): its coefficients are integer multiples of powers
of $\Gs$, its roots $x_+=137.036\ldots$ and $x_-=3.024\ldots$ are an exercise in
the quadratic formula, and the whole of it is a theorem. To turn that polynomial
toward physics, one more thing is needed, and naming it sharply is the entire
content of this Part.

\runin{The readout move.} Read the quadratic as the characteristic polynomial of
a $2\times2$ \emph{readout operator} $W$, an admissible operational
functional on the substrate whose dominant eigenvalue is to be the measured
inverse coupling. For its characteristic polynomial to be $P$, the operator must
have
\[ \assembly. \]
This is the \textbf{operator assembly}. The question that decides whether
$\alpha$ is derived or merely matched is exactly:

\begin{boundarybox}[The readout problem \etag{opn}]
Does $\Ftheory$ (the five postulates plus the algebraic spine plus the
$O_h$ representation theory of the BCC corner module) \emph{force} the
assembly $(\Tr,\Det)=(16\Gss,16\Gsss)$, or is that assembly an \emph{imposed
selection}?
\end{boundarybox}

\noindent This is the readout-structure criterion, and the contract a closure
must satisfy is fixed in advance and deliberately strict: an admissible readout
must be stated \emph{before} any physical value is checked, and a closure fails
immediately if it inserts physical $\alpha$, a CODATA value, $g_c$, or $x_+$, or
uses the framework's transfer matrix $M_N(t)$, which is \emph{defined} to
have the master quadratic as its characteristic polynomial, hence circular and
banned. The missing object is not another fit; it is an operational readout rule.
Part~II determines whether the ontology supplies one.

\runin{Why this is the sharp form of the obstruction.} The central foundational
obstruction of the framework's $\alpha$ program comes down to this one structural
question, because everything \emph{around} it is already settled in Part~I. The
constant $\Gs$ is forced (four ways, \secref{sec:gstar}); the coefficient
$16=|\Ziu|^2$ is a theorem at the value level (\secref{sec:masterquad}); the
polynomial's roots are a theorem. If the \emph{assembly} were forced too, $\alpha$
would be \etag{der}. So the entire weight of the derive-or-match verdict rests on
one hinge: whether $\Ftheory$ produces $W$.

\section{The four FTD-native routes}\label{sec:routes}
A boundary cannot be argued from absence, from merely failing to find an
assembly. The canonical attack instead ran \textbf{four FTD-native routes of
differing character} (later shown, in \secref{sec:routeinv}, to be four dialects
of one underlying obligation, so ``0 of 4'' is one structural fact reached four
ways, not four independent failures). Each was first \emph{force-attempted} (build
the strongest honest forward chain to the assembly, $\Gs$ kept symbolic, all
banned moves excluded) and then \emph{adversarially refuted} by a second pass
hunting for any smuggled step:

\begin{table}[h]
  \centering\small
  \renewcommand{\arraystretch}{1.3}
  \begin{tabularx}{\linewidth}{@{}l >{\RaggedRight\arraybackslash}X@{}}
    \toprule
    \textsc{route} & \textsc{FTD-native principle attempted}\\
    \midrule
    \texttt{jtwist} & the $J$-twisted $\zeta$-regularized determinant operator
      (the candidate clean odd source)\\
    \texttt{bcc} & the BCC body-diagonal transfer / response operator
      (triple-cosine Watson)\\
    \texttt{cm} & the lemniscatic-CM arithmetic of $E:y^2=x^3-x$
      (CM by $\Zi$, $\Aut=\mu_4$)\\
    \texttt{novel} & a forced variational principle / period-ring (Hodge)
      valuation / $K$-theory co-realizability\\
    \bottomrule
  \end{tabularx}
\end{table}

\runin{The verdict: 0 of 4 force the assembly.}\margetag[ no-go]{smc} Every
force-pass self-reported a \texttt{gap}; every refute-pass upheld
\texttt{boundary}; the set of cleanly-forced routes is empty. The honest reading
partitions cleanly into what the ontology \emph{does} force and what it does
\emph{not}:
\begin{itemize}[leftmargin=1.4em, itemsep=4pt]
  \item \textbf{What IS forward-forced} (\etag{der}), agreed by every route with
        nothing smuggled:
    \begin{itemize}[leftmargin=1.2em, itemsep=2pt]
      \item the \textbf{even trace} $16\Gss$: $16=|\mu_4|^2$ for
            $E:y^2=x^3-x$; $\Gss=2\pi\cdot G_{\mathrm{BCC}}(0)$, the Watson BCC
            self-energy;
      \item the \textbf{existence of a clean FTD-native odd source}: the
            $J$-twisted determinant ratio
            $\det_\zeta(D_{3/4})/\det_\zeta(D_{1/4})=\Gamma(1/4)/\Gamma(3/4)=\Gs$
            (degree~1; the $\sqrt{2\pi}$ cancels, so there is no forbidden
            $\sqrt\pi$ prefactor). This is \emph{strictly stronger than
            trace-only}: it means $16\Gsss=16\Gss\cdot\Gs$ is
            \textbf{assemblable}, genuinely lifting the bare-parity no-go.
    \end{itemize}
  \item \textbf{What is NOT forced} (\etag{opn}):
    \begin{itemize}[leftmargin=1.2em, itemsep=2pt]
      \item the \textbf{operator assembly itself}: that the \emph{same}
            $2\times2$ readout carrying $\Tr=16\Gss$ also has $\det=16\Gsss$. For
            a $2\times2$ operator, trace and determinant are \textbf{independent
            invariants}; fixing the trace leaves the determinant free. The odd
            scalar $\Gs$ exists, but nothing forces the \emph{gluing}: neither
            that the two invariants belong to one operator, nor that the odd
            scalar lands in the determinant slot rather than anywhere else. This
            residual is precisely the imposed Vieta target of the
            readout-structure criterion.
    \end{itemize}
\end{itemize}
The differing character of the four channels is what makes ``0/4'' load-bearing
rather than anecdotal; and \secref{sec:routeinv} turns the fact that they
reach one obligation into a theorem. The overall no-go is a \etag{smc}, not a
theorem: the routes are adversarial refutation \emph{attempts}; none constructed
the forcing assembly, and none proved one \emph{cannot} exist. Calling the
boundary a theorem on this evidence would overreach. It is a boundary-mapping
deliverable in the exact sense of Clause~2, at honest conjecture grade.

\runin{A representative route, made concrete.} The boundary shows itself cleanly
in the modular dialect. At the self-dual point $\tau=i$ (the fixed point of the
modular $S$-transformation $\tau\mapsto-1/\tau$), the torus partition function's
modular geometry forces $E_6(i)=0$ and $E_4(i)\propto G^{*4}$, so the available
invariants are \emph{even} powers of the period $\Gs$ ($\Gss$ from the
Green's-function variance, $G^{*4}$ from $E_4$). The required determinant
$16\Gsss$ is an \emph{odd} power, which no theta or modular form on the torus
generates without inserting an external scale. The verdict is exactly
\texttt{UNDERDETERMINED}: this route supplies the even trace and cannot force the
odd determinant slot, one instance of the route-invariant pattern the next
section names.

\section{Route-invariance and the theorem-grade components}\label{sec:routeinv}
The deeper attack asked whether the residual gap could be closed by a sharper
instrument; and in failing to close it, established three genuinely
theorem-grade facts plus the reason all four routes land in the same place. These
are the boundary's \emph{masonry}: each is a theorem, and together they hem the
open kernel of \secref{sec:kbind} into a single, precisely-located gap.

\runin{(a) The flip is ruled out.}\margetag{thm} A \emph{flip} would be a
substrate-native operator that genuinely \emph{forces} $\det=16\Gsss$, flipping
the verdict to ``$\alpha$ derived.'' The only geometric candidate is the
$C_3(\langle111\rangle)$ three-plane $\det_\zeta$ product (a $D_6$-symmetric
object built from three cyclically-permuted planes, each contributing one factor
of $\Gs$). It is \textbf{excluded} from the rank-2 readout in two steps:
\begin{itemize}[leftmargin=1.4em, itemsep=2pt]
  \item a definite complex structure $J$ with $J^2=-I$ (required by the
        trace via the BCC $V_{\mathrm{complex}}=\Zi^2$ structure) needs
        $\mult_O(E)=0$ (the character table of $O$ shows no $O$-symmetric
        2-dimensional subspace of the 8-corner module), which forces $O_h$ to
        break to \textbf{one $C_4$ axis}, so $C_3(\langle111\rangle)\notin\mathrm{Stab}$;
  \item the $D_6$ three-plane product requires
        $C_3(\langle111\rangle)\in\mathrm{Stab}$;
  \item but a $4$-fold face axis and a $3$-fold body-diagonal axis generate the
        whole octahedral rotation group, $\langle C_4, C_3\rangle = O$ (order
        24), so demanding both restores unbroken $O_h$ $\Rightarrow$ no localized
        charge $\Rightarrow$ no $V_{\mathrm{complex}}$ $\Rightarrow$ no readout.
\end{itemize}
The two requirements are mutually exclusive from a single preparation. Every
$2\times2$ that \emph{does} realize $(16\Gss,16\Gsss)$ (e.g.\ the companion
form $\bigl[\begin{smallmatrix}0 & -16\Gsss\\ 1 & 16\Gss\end{smallmatrix}\bigr]$) has
its determinant entry \emph{hand-placed}, which is exactly the banned gluing: a
witness for the cheap argument, not a derivation.

\runin{(b) The equivariant restriction closes its own scope.}\margetag{thm} No
$C_3(\langle111\rangle)$-equivariant rank-2 \emph{restriction} of the three-plane
source carries $(\Tr,\Det)=(16\Gss,16\Gsss)$. On the $C_3$-adapted basis
$\R^3=R\oplus C$, the unique rank-2 $C_3$-invariant subspace is the $C$-plane,
whose $C_3$-commutant is $\{xI+yK\}\cong\C$ (Schur). The master-quadratic roots
are real and distinct ($\mathrm{disc}=64\Gsss(4\Gs-1)>0$), and a
$C_3$-equivariant $C$-plane operator has real-distinct spectrum \emph{iff it is
non-real}, which invokes the ambient scalar $i$, supplied only by breaking
$O_h$ to one $C_4$ axis, whence $\langle C_4,C_3\rangle=O$ and there is no
readout. The chain is therefore \textbf{reality $\Rightarrow$ scalar-$i$
$\Rightarrow$ $C_4$ $\Rightarrow$ $O$}: the \emph{same} wall as the rest of
the program, reached from the restriction side, not an independent conjugacy
fact. The same mechanism holds dimension-free: any $C_3$-equivariant linear
reduction of the three-plane object factors through the $C_3$-fixed diagonal
(where $C_3$ acts as the identity), collapsing the determinant grading
$\Gsss\to G^{*1}$; the rank-2 compressions of $\Gs\cdot I_3$ have determinant
$\Gss$ at most, never $\Gsss$. This closes the restriction's scope; the surviving
non-$C_3$-invariant and det-by-fiat cases are exactly the open kernel.

\runin{(c) The reduction is route-invariant.}\margetag{thm} All four attacks
reduce to the \emph{same} obligation in four dialects, and this is not a
coincidence; it follows from one clean field-theoretic fact:
\begin{quote}\itshape
  $\Q(\Gs)$ is the Galois-fixed field of the master quadratic's $\Z/2$.
\end{quote}
\noindent The quadratic's discriminant is $\Delta=64\Gsss(4\Gs-1)$, so its
splitting field over $\Q(\Gs)$ is $\splitfield$, with Galois group $\Z/2$ swapping
$x_+\leftrightarrow x_-$. Every forward-forced \emph{symmetric} FTD-native
datum (the Watson trace $16\Gss$, the $\det_\zeta$ ratio $\Gs$, the
Chowla--Selberg periods) lives in the $\sigma$-fixed subfield $\Q(\Gs)$ and
is therefore \textbf{provably blind to which root is $1/\alpha$}. The concrete
witness: the family $\det=16\Gss\cdot\Gpow{k}$ for $k=0,1,2,3$ gives dominant
roots $139.05\,/\,137.04\,/\,130.68\,/\,105.76$, \emph{all} of them
$\Ftheory$-consistent, and \textbf{nothing in $\Ftheory$ selects $k=1$.} The
regularized traces $\zeta(-1,1/4)=\zeta(-1,3/4)=1/96$ are rational, so the trace
carries zero $\Gs$; and at rank~2 the regularized determinant is the ordinary
product $x_+\cdot x_-=16\Gsss$, so hitting the assembly is a \emph{chosen}
Vieta target, not a forced one.

\medskip
\noindent These three components are why the boundary is a \emph{wall} and not a
mere absence of progress: the obvious forcing candidate is dead~(a), the
equivariant restriction is closed within its scope~(b), and any remaining route is
provably the \emph{same} route in disguise~(c). What is left is one gap, and
\secref{sec:kbind} states it exactly.

\section{The square-root wall}\label{sec:kbind}
The three theorem-grade components hem the obstruction into a single irreducible
obligation. Stated once, plainly:

\begin{boundarybox}[The binding obligation \etag{opn}]
Prove (or refute) that no substrate-native operator construction can
bind the degree-1, $C_3$-agnostic odd scalar $\Gs$ (the $J$-twisted $\det_\zeta$
ratio) into the determinant slot of the \emph{same} $2\times2$ readout that
carries the definite complex structure $i$, with the exponent fixed at exactly
\textbf{1} by the substrate rather than chosen.
\end{boundarybox}

\noindent The basis-free fact is the primary one: trace and determinant are
independent invariants of a rank-2 operator over $\Q(\Gs)$, so the assembly $W$ is
a free choice. The square-root form below is the cleanest single-coordinate
dialect of that one fact, a basis-representative of a basis-free obstruction,
not a coordinate-bound claim (the determinant's parity is itself basis-dependent:
a rescaling migrates the odd content into the trace). In that dialect: to fix
$\alpha$, the substrate must \textbf{natively realize}
\[ \sqrtwall, \]
the \textbf{squarefree generator} of the quadratic extension
$\Q(\Gs)(\sqrt\Delta)/\Q(\Gs)$; since $\Delta=64\Gsss(4\Gs-1)$ we have
$\sqrt\Delta=8\Gs\cdot\sqrtwall$, the irrational quantity that \emph{first
distinguishes} $x_+$ from $x_-$. Every \emph{symmetric} datum lives in $\Q(\Gs)$
and cannot tell the two roots apart (\secref{sec:routeinv}(c)); the square root is
exactly the ingredient that breaks the $\Z/2$ symmetry and picks a root. This is
the discrete-limit-of-the-square-root wall: the precise place the construction
halts. The substrate forces every \emph{symmetric} ingredient of the quadratic;
it does not, on present evidence, produce the \emph{one antisymmetric
beable} (the square root) that would select the physical branch.

\runin{Why it is \etag{opn}, and not closeable today.} The obligation is a
\textbf{universal negative} over substrate-native operator constructions, and it
is not yet a well-posed proposition, because $\Ftheory$ does not yet pin down the
class it quantifies over. Provisionally, take that class to be the $2\times2$
operators on $V_{\mathrm{complex}}=\Zi^2$ generated by the substrate-native
data (the BCC corner module, the $C_4$-winding preparation, and the
$\det_\zeta$ functor) closed under direct sum and restriction to a stabilizer
(the regularized determinant furnishing the characteristic invariants, not itself
a closure operation). Whether this is the \emph{right} closure, whether a
larger one (tensor products, non-self-adjoint generators) would change the
verdict, is itself part of the open work, not settled here. Until that calculus is fixed precisely, ``no
substrate-native operator binds the odd scalar at exponent~1'' quantifies over an
undefined domain, and one cannot prove a universal over a class one has not
finitely specified. Two senses of \etag{opn} are worth separating: an open
\emph{conjecture} is a well-posed proposition awaiting proof; an open
\emph{desideratum} is a property not yet expressible as a proposition because its
quantifier-domain is unspecified. The binding obligation is currently the latter;
fixing the operator calculus would convert it to the former. \textbf{Crucially,
the conditional theorem of \secref{sec:conditional} is insulated from this}: it
rests only on trace and determinant being independent rank-2 invariants over
$\Q(\Gs)$ (which is well-posed and proved), so the boundary's logic does
not wait on the operator calculus.

\runin{The one non-axiomatic exit is currently shut.} The alternative to
axiomatizing the calculus is to \emph{measure} the binding directly on the
substrate. That measurement has been run, and it returned a \etag{cn}: a null
robust to measurement precision, \textbf{0 macroscopic cluster fissions
across 2000 seeds} in the tested regime, from a protocol that would have detected
fissions above its sensitivity floor and saw none. One should not read
``impossible'' into a finite null; but neither route, the axiomatic (no
calculus) nor the empirical (a null result so far), is presently available, and
this is the single remaining obligation of the entire $\alpha$ program.

\section{The conditional theorem}\label{sec:conditional}
What \emph{is} provable, rigorously and unconditionally, is the conditional. It is
the strongest theorem-grade statement the boundary supports, and the precise sense
in which ``$\alpha$ is not forced'' is earned rather than asserted:

\begin{forcedbox}[The conditional theorem \etag{thm}]
Nothing that $\Ftheory$ forces selects the operator assembly
$(\Tr,\Det)=(16\Gss,16\Gsss)$. Every datum the spine forward-forces is
\emph{symmetric}; it lies in the fixed field $\Q(\Gs)$ of the master
quadratic's $\Z/2$, and is therefore blind to which root is $1/\alpha$. With the
trace $16\Gss$ already forward-forced, only the determinant slot is free: the
assembly is one of an infinite family of determinant values over that fixed
trace, all equally compatible with everything $\Ftheory$ forces. To single it out (equivalently,
to realize a beable in $\Q(\Gs)(\sqrtwall)\setminus\Q(\Gs)$) demands a
substrate-native binding law $W$ beyond $\Ftheory$ that fixes the determinant's
odd-$\Gs$ exponent at exactly \textbf{1} from a single stabilizer.
\end{forcedbox}

\noindent The force of this conditional is carried by the \secref{sec:routeinv}
components, above all the Galois-fixed-field fact, which is what makes
``$\Ftheory$ does not select the assembly'' a theorem and not a restatement of
ignorance: the two sides are exhibited explicitly, as realizations, not merely
asserted.
\begin{itemize}[leftmargin=1.4em, itemsep=3pt]
  \item \textbf{Adopting $W$ is coherent}, the master quadratic is the
        explicit realization: once the assembly is posited it is internally
        consistent.
  \item \textbf{Declining $W$ is equally coherent}: $\det=16G^{*4}$ (dominant
        root $\approx130.68$) and $\det=\Gs$ (dominant root $\approx140.04$) are
        explicit alternative assemblies that everything the spine forces equally
        permits.
\end{itemize}
Because both are coherent, no binding law is implied by what $\Ftheory$
forces; and \emph{that}, not any failed search, is the content of ``the
ontology does not force the assembly.'' $W$ is, in the framework's bookkeeping, a
\textbf{6th-postulate-class input}: added to FTD's axioms \emph{and} shown
substrate-native, it would convert $x_+=1/\alpha$ from a \etag{smc} to
\etag[, modulo $W$]{der}. Whether any FTD-native $W$ exists is exactly the binding
obligation of \secref{sec:kbind}, which remains \etag{opn}. The conditional
itself is theorem-grade; the \emph{unconditional} no-go (``no FTD-native $W$
can ever exist'') is the universal negative, which stays open, and is why the
overall boundary is a \etag{smc}.

\medskip
\noindent This is the exact meaning of the shorthand used elsewhere in the
monograph. ``$\Ftheory$ does not force $\alpha$'' means: no selection follows
from the present five-postulate theory plus the spine; a future $W$ would be a
new binding law, not a consequence already hidden in the current axioms. It does
\emph{not} mean that no possible extension of FTD could ever force the coupling.

\section{Conclusion: \texorpdfstring{$\alpha$}{alpha} is dynamical, not structural}\label{sec:dynamical}
The deliverable of Part~II is a \textbf{mapped wall}. In the project's own
phrasing: \textbf{the discrete ontology forces the \emph{menu}} (the even
trace $16\Gss$ from the Watson integral, and the existence of a clean odd source
$\Gs$ from the $J$-twisted $\zeta$-determinant), \textbf{but it does not force
the \emph{dish}}, the assembly of those ingredients into one readout operator $W$.
So $\alpha$'s value rides on an independent convention, not on the postulates.
The assembly is unforced on present evidence: that contrast is \etag{der};
the reading it licenses, \textbf{$\alpha$ is dynamical, not structural}, is the
\etag{smc} that summarizes it, not a proof that no future binding law exists. The
contrast with $N_c$ is exact and
instructive: $N_c=3$ \emph{is} structural, forced from $O_h$ and topology by
four independent routes with no operator-assembly choice required; the value falls
out of the symmetry. $\alpha$ has no such forcing: with the trace forward-forced,
everything the ontology fixes is still consistent with infinitely many determinant
values over that one trace, the same epistemic status the engine assigns its coupling
$g_c$.

\medskip
\noindent Accordingly, the central obstruction remains a \etag{fo}, and
$x_+=1/\alpha$ remains a \etag{smc}; nothing in this Part promotes either. The
wall is built of theorem-grade masonry (the ruled-out flip, the closed-scope
restriction, the route-invariant reduction, and the conditional theorem) around
one clearly-marked open gap: the square root that would break the root-swapping
symmetry. The surviving exits are explicit and few: a 6th-postulate-class input
$W$ that \emph{forces} the assembly, or a fresh discrete-native measurement (and
the one run so far returned a \etag{cn}).

\medskip
\noindent This is exactly what Clause~2 of the governing goal asks for.
\emph{``Rigorously establish the boundary of what the ontology cannot force''} is
not a consolation prize for failing to derive $\alpha$; it is a deliverable in
its own right. Part~II draws that line where the construction actually halts: at
the discrete limit of a square root. A program that only exhibited the spine of
Part~I would be a curiosity; one that maps this wall, in both directions and at
honest grades, is doing the thing the goal was set to do.
